(a) $\bar{\mathbf{X}}$

(b) $n(\bar{\mathbf{X}}-\boldsymbol{\mu})'\mathbf{S}^{-1}(\bar{\mathbf{X}}-\boldsymbol{\mu})$

4.23\quad 考虑 1996—2005 年间道琼斯工业股票价格指数的年回报率 (包括红利). 这些数据乘以 100, 便是 $-0.6,3.1,25.3,-16.8,-7.1,-6.2,25.2,22.6$ 和 $26.0$. 试利用这 9 个观测值完成下列问题.

(a) 作一个 Q-Q 图. 这些数据看起来是否正态分布? 请解释.

(b) 完成一个基于相关系数 $r_{Q}$ 的正态性检验 [见\hyperref[eq:4-31]{式(4-31)}]. 显著性水平取 $\alpha=0.10$.

4.24\quad 练习 1.4 给出了 2005 年 4 月 10 家最大公司三个变量的数据. 试对销售额 $(x_{1})$ 和利润 $(x_{2})$ 的数据

(a) 作 Q-Q 图. 这些数据看起来是否正态分布? 请解释.

(b) 完成一个基于相关系数 $r_{Q}$ 的正态性检验 [见\hyperref[eq:4-31]{式(4-31)}]. 显著性水平取 $\alpha=0.10$. 这些检验结果能否证实 (a) 中的结果?

4.25\quad 参看练习 1.4 中 10 家最大公司的数据. 利用所有三个变量作卡方图. 卡方分位数为 $0.3518$ $0.7978$ $1.2125$ $1.6416$ $2.1095$ $2.6430$ $3.2831$ $4.1083$ $5.3170$ $7.8147$

4.26\quad 练习 1.2 给出 $n=10$ 辆旧车的车龄 $x_{1}$ (年) 和销售价格 $x_{2}$ (千美元). 这些数据如下:

\begin{mvtable}
\begin{tabular}{lrrrrrrrrrr}
\toprule
$x_{1}$&1&2&3&4&5&6&7&8&9&11\\
\midrule
$x_{2}$&18.95&19.00&17.95&15.54&14.00&12.95&8.94&7.49&6.00&3.99\\
\bottomrule
\end{tabular}
\end{mvtable}

(a) 利用练习 1.2 的结果计算平方统计距离
\[
(\mathbf{x}_{j}-\bar{\mathbf{x}})'\mathbf{S}^{-1}(\mathbf{x}_{j}-\bar{\mathbf{x}}),j=1,2,\cdots,10,\quad\text{其中 } \mathbf{x}_{j}'=[x_{j1},x_{j2}].
\]

(b) 利用 (a) 中得出的距离, 确定落在二维正态分布估计的 $50\%$ 概率轮廓线之内的诸观测值所占的比例.

(c) 将 (a) 中的距离排序并作卡方图.

(d) 给定 (b) 和 (c) 的结果, 这些数据是否接近二维正态分布? 请解释.

4.27\quad 考虑例 4.10 的辐射数据 (炉门关闭). 对这些数据的自然对数作 Q-Q 图 [注意自然对数变换对应\hyperref[eq:4-34]{式(4-34)} 中的 $\lambda=0$]. 这些自然对数看起来是否呈正态分布? 将你的结果同\hyperref[fig:4.13]{图 4.13} 相比较. 在这种情况下, 选择 $\lambda=1/4$ 或 $\lambda=0$ 是否会产生很大差别?

\textbf{下面的练习需要使用计算机.}

4.28\quad 考虑\hyperref[tab:1.5]{表 1.5} 中所给的空气污染数据, 对太阳辐射测量值作 Q-Q 图, 并根据相关系数 $r_{Q}$ 完成正态性检验 (见\hyperref[eq:4-31]{式(4-31)}). 取 $\alpha=0.05$ 并利用\hyperref[tab:4.2]{表 4.2} 中对应于 $n=40$ 的 $r_{Q}$ 值.

4.29\quad 用\hyperref[tab:1.5]{表 1.5} 中所给的空气污染数据, 考察 $X_{5}=\mathrm{NO}_{2}$ 和 $X_{6}=\mathrm{O}_{3}$ 这对变量的二维正态性.

(a) 计算统计距离 $(\mathbf{x}_{j}-\bar{\mathbf{x}})'\mathbf{S}^{-1}(\mathbf{x}_{j}-\bar{\mathbf{x}}),j=1,2,\cdots,42$, 其中 $\mathbf{x}_{j}'=[x_{j5},x_{j6}]$.

(b) 确定落在二维正态分布的大约 $50\%$ 概率轮廓线之内的观测值 $\mathbf{x}_{j}'=[x_{j5},x_{j6}],j=1,2,\cdots,42$, 所占的比例.

(c) 用 (a) 中排序距离作卡方图.

4.30\quad 考虑练习 4.26 中旧车数据.

(a) 确定能使 $x_{1}$ 值接近正态分布的幂变换 $\hat{\lambda}_{1}$. 对变换后的数据作 Q-Q 图.
